% LSPV-4 REPORT::ORDERBIND1::2026-10-12T00:01:48Z
\documentclass[11pt]{article}
\usepackage[a4paper,margin=0.86in]{geometry}
\usepackage{amsmath,amssymb,amsthm,booktabs,array,enumitem,xurl,microtype}
\usepackage[hidelinks]{hyperref}
\setlength{\emergencystretch}{3em}
\newcommand{\SAT}{\mathsf{SAT}}
\newcommand{\EXACT}{\textbf{PROVED/EXACT}}
\newcommand{\KNOWN}{\textbf{KNOWN THEOREM}}
\newcommand{\COMP}{\textbf{COMPUTATIONAL}}
\newcommand{\CAND}{\textbf{CANDIDATE}}
\newcommand{\OPEN}{\textbf{OPEN}}
\newcommand{\RETIRED}{\textbf{FALSIFIED/RETIRED}}
\newtheorem{theorem}{Theorem}
\newtheorem{corollary}{Corollary}
\title{LSPV-4 ORDERBIND1: Exact source-order binding costs only permutation entropy}
\author{UNIVERSAL\_CORE\_BRIDGE / LSPV-4}
\date{12 October 2026, 00:01:48 UTC}

\begin{document}
\maketitle

\begin{abstract}
The current synthesis asks whether a source-matched BIND theorem can repair the
collapse of a physical locator under equivalent clause orders.  This report
separates two meanings of BIND.  For a duplicate-free $q$-clause formula, an
explicit tag that reconstructs the \emph{original ordered encoding} from the
unordered clause multiset must distinguish all $q!$ orderings and therefore
has at least $\lceil\log_2(q!)\rceil$ bits.  This bound is tight by Lehmer
coding and is only $\Theta(N)$ on the active exact-$(2,2)$ encoding.  If the
decoder needs only the clause multiset or a canonical representative, the
entire permutation is quotientable and no source-order tag is required.
Thus order binding is either representation-specific linear output or
semantically unnecessary; it does not force compatibility extraction, a
hereditary pass, or difficult one-bit SAT decision.  The surviving target must
bind a semantic child/opening needed by exact reconstruction, not a harmless
input permutation.  P versus NP remains \OPEN.
\end{abstract}

\section{Claim card and exact quantifiers}

\begin{center}\scriptsize
\begin{tabular}{>{\raggedright\arraybackslash}p{0.32\linewidth}>{\raggedright\arraybackslash}p{0.22\linewidth}>{\raggedright\arraybackslash}p{0.38\linewidth}}
\toprule
Claim & Status & Scope\\ \midrule
\texttt{ORDER-SOURCE-BINDING-DICHOTOMY1} & \EXACT, independent review pending & Fixed-length explicit tags over a duplicate-free clause-permutation orbit.\\
Finite orbit audit & \COMP & $2{,}863{,}907$ checks through $q=9$; zero failures.\\
\texttt{ORDER-BINDING-AMPLIFIES-HARDNESS1} & \RETIRED & Exact source order yields only a tight linear representation/output charge; canonical semantics quotients it.\\
\texttt{SEMANTIC-OPENING-BIND1} & \OPEN & Two conflicting valid openings must yield a construction-specific forbidden witness.\\
Arbitrary-decider transfer and P versus NP & \OPEN & No model-matched contradiction.\\
\bottomrule
\end{tabular}
\end{center}

Let $M=\{C_1,\ldots,C_q\}$ be a set of $q$ pairwise distinct fixed-width
clause records and let
\[
 \operatorname{Orb}(M)=\{(C_{\pi(1)},\ldots,C_{\pi(q)}):\pi\in S_q\}.
\]
The attacked repair is that retaining a source-covariant certificate for this
orbit creates a nontrivial universal SAT obligation.  Its exact reconstruction
version has quantifier prefix
\begin{equation}
 \exists b\;\exists T:\operatorname{Orb}(M)\to\{0,1\}^b\;
 \exists D\;\forall E\in\operatorname{Orb}(M):
 D(M,T(E))=E.                                                   \tag{1}
\end{equation}
The SAME\_REALIZATION objects are the immutable source encoding $E$, the
emitted or stored tag $T(E)$, and the decoder $D$.  Construction, output,
storage, access and replay of the tag are charged.  No SAT decider appears in
(1); promoting (1) to arbitrary decision computation is a separate \OPEN
transfer.

\section{Exact orbit theorem}

\begin{theorem}[\texttt{ORDER-SOURCE-BINDING-DICHOTOMY1}]
For every $M$ as above:
\begin{enumerate}[leftmargin=1.7em]
\item If (1) holds, then $T$ is injective and
\[
 b\ge \left\lceil\log_2(q!)\right\rceil.
\]
This is tight by factorial-number-system (Lehmer) coding.
\item For an arbitrary tag $T$ with fibers $f_z=|T^{-1}(z)|$ and collision
count $\operatorname{Coll}(T)=\sum_z f_z^2$,
\[
 |\operatorname{im}T|\ge
 \frac{(q!)^2}{\operatorname{Coll}(T)}.                        \tag{2}
\]
Exact reconstruction is equivalent to $\operatorname{Coll}(T)=q!$.
\item If the required output is only $M$ or any fixed canonical ordering of
$M$, then the constant tag suffices; its single fiber has size $q!$ and
collision count $(q!)^2$.
\end{enumerate}
\end{theorem}

\begin{proof}
If $T(E)=T(E')$, applying $D(M,\cdot)$ in (1) gives $E=E'$.  Hence $T$ is
injective, so $2^b\ge|\operatorname{Orb}(M)|=q!$.  Conversely the Lehmer rank
is a bijection between $S_q$ and $\{0,\ldots,q!-1\}$ and uses the asserted
number of bits.  Equation (2) is Cauchy--Schwarz:
$(\sum_z f_z)^2\le |\operatorname{im}T|\sum_z f_z^2$.  Injectivity makes
every nonempty fiber one; canonicalization maps the entire orbit to one
representative.  \qed
\end{proof}

On proper Monotone $3$SAT-$(2,2)$ with $n$ variables there are
$q=4n/3$ clauses.  The canonical indexed encoding has
$N=\Theta(q\log n)=\Theta(n\log n)$, while Stirling's formula gives
\begin{equation}
 \log_2(q!)=q\log_2 q-(\log_2 e)q+O(\log q)=\Theta(N).          \tag{3}
\end{equation}
Thus exact physical source-order recovery gives a tight linear-bit output or
representation obligation, not a superlinear or superpolynomial SAT lower
bound.  If the project fixes a canonical clause order, the permutation orbit
is absent from the input contract; if it accepts arbitrary order but treats
CNF as a conjunction, semantic decision may quotient the orbit.

\begin{corollary}[No locator-to-hardness repair]
Appending exact source-order BIND to ORDERLOC1 does not turn its locator floor
into an arbitrary-decider lower bound.  The tag can be a tight permutation
code, while the one-bit SAT value is constant over the orbit and need not
retain any order information.  Therefore
\texttt{ORDER-BINDING-AMPLIFIES-HARDNESS1} is \RETIRED.
\end{corollary}

This does not refute every possible BIND theorem.  It removes physical order
as the object.  A surviving construction must define an opening whose
conflicting valid values alter a hereditary semantic child or produce an
explicit forbidden algebraic witness.  That is
\texttt{SEMANTIC-OPENING-BIND1}, and it remains \OPEN.

\section{Bounded reproducible audit --- \COMP}

The dependency-free program
\path{experiments/review/order_binding_orbit_audit.py} exhaustively enumerates
all $q!$ permutations for $1\le q\le9$.  It verifies Lehmer
encoding/decoding, the exact bit ceiling, (2), and the fibers for first record,
last record, two endpoints, inversion parity and canonical sorting.  At
$q=9$, the orbit has $362{,}880$ elements: the exact tag has $362{,}880$
singleton fibers and needs 19 bits, whereas canonical sorting has one fiber
and collision count $131{,}681{,}894{,}400=(9!)^2$.

The same bounded invocation checks DIRECT, COPY, DEAD, DELAY, INTERP and
RECOMP one-bit wrappers on every orbit element.  It made $2{,}863{,}907$
assertions with zero failures in 3.3281 seconds.  The raw result is
\path{results/raw/review/2026-10-12_order_binding_orbit.json}; its canonical
payload SHA-256 is
\begin{center}
\texttt{ebc8b17a7cb576580ac2e8e2b80a3c2511de7b6ca1066c4c70f7a0a0f3db5b39}.
\end{center}
The finite audit supports implementation correctness only.  The asymptotic
claim is the symbolic theorem above.  The single run is exhaustive at its
declared finite sizes and remains well below the repository's 30-second
limit; no chunking was needed.

\section{Hostile controls and SAME\_REALIZATION}

Clause permutation preserves satisfiability and membership in Horn, dual
Horn, 2SAT, affine/Tseitin, bounded-treewidth and proper bounded-occurrence
families.  Independent products multiply harmless ordering choices.  Padding
or duplicate inert records can inflate permutation entropy without making the
one-bit decision harder, so they are hostile rather than supporting examples.
Variable renaming and complementation add further representation orbits.
Canonical sorting, characteristic-product equality, succinct permutation
codes, indexes and compact circuits are alternative representations that
must be charged when actually used but cannot be declared mandatory for an
arbitrary SAT execution.

If $T(E)$ is explicitly emitted, (3) is an output-length lower bound for that
chosen transducer.  If it is stored for random access, it is a data-structure
space obligation.  If it is recomputed, its computation is charged.  None of
these cases proves that a one-bit decider instantiates $T$, and mixing the
tag-producing machine with the original decider would violate
SAME\_REALIZATION.  A maximal semantic SAT oracle discards order entirely;
an algebraic oracle can do the same.  The theorem remains true but irrelevant
to their decision cost.

\section{Literature and cross-domain transfer audit}

Munro--Raman--Raman--Rao study the exact same permutation representation
object and give succinct encodings near the information-theoretic
$\lceil\log_2(q!)\rceil$ floor with operation-time tradeoffs
\cite{MRRR12}.  This is a \textbf{directly transferable theorem} only for an
explicit permutation data structure: records correspond to clauses and the
stored permutation corresponds to $T(E)$.  The missing SAT theorem is that an
arbitrary decider must support those permutation operations or reconstruct
the source order.  No such theorem is present.

FLMV's polynomial SAT time--space bounds apply after one ordinary random-
access machine is placed in their simultaneous resource regime
\cite{FLMV05}.  A linear tag-output floor does not do so.  VLSI routing,
communication cuts, distributed consensus and control observability are
\textbf{misleading analogies / failed transfers} here: clause positions are
not fixed processors, terminals or hidden states.  Canonical labeling and
succinct indexing are \textbf{useful structural analogies requiring a new
transfer theorem}; they solve representation access, not SAT inevitability.

The post-quantum collision ledger is an \textbf{EXACT OBJECT/THEOREM
TRANSFER}: (2) is the prescribed lossiness bound.  Schur squares, quadratic
relation ranks, Hilbert functions and syzygies are \textbf{FAILED TRANSFER}
because no code is forced.  Special-soundness/collision extraction is an
\textbf{ANALOGY/CANDIDATE}: two different clause-order openings are both
valid and yield no forbidden witness, so a future construction needs a
semantic inconsistency theorem rather than generic uniqueness.  Hashes,
random oracles and cryptographic binding assumptions remain inadmissible.

The counting proof relativizes, and oracle access makes the missing transfer
strictly easier to bypass; Baker--Gill--Solovay therefore remains applicable
\cite{BGS75}.  Nothing resists low-degree extension access, so the argument
does not cross algebrization \cite{AW09}.  It constructs no useful large
circuit property, so Natural Proofs is neither used nor evaded.

\section{CORPUS REVIEWED}

The live GitHub tree was enumerated before substantive work.  This run read
root \path{main.tex}; actual \path{main4.tex}/\path{main5.tex};
\path{AGENTS.md}; \path{CODEX.md}; the complete pnp-research skill, mode and
research-protocol files; and all current task specifications under
\path{tasks/}.  The newest verified historical MAIN was
\path{research/main/2026-10-11_2348_MAIN_SYNTHESIS-0017.tex}.  Material peer
reports included LSPV-1 LB-0020, LSPV-2 TESTS-0016, LSPV-3 ORDERLOC1, LSPV-4
TAGGLUE1, MODELBRIDGE-0015/0016/0019 and SYNTHESIS-0017.  Relevant older
source-order work in ARC10--ARC15, especially ARC13's characteristic-product
multiset comparison, was inspected rather than silently duplicated.

The HEC, VSC-HEC, G-PPC-HEC, CRS-PPC, CCSPC and RCCN causal branch, governance
notes, coordination/spine files, and cited experiment/result artifacts were
reviewed.  Repository search covered clause order, permutation orbits,
canonicalization, source-matched BIND, collision entropy, locator and
SAME\_REALIZATION.  Primary literature \cite{MRRR12,FLMV05,BGS75,AW09} was
checked before promotion.

The cyclic predecessor
\path{research/agents/lspv-3/2026-10-11_2334_LSPV-3_ORDERLOC1.tex} exists at
its canonical GitHub path, is substantive, contains the required header and
final sentinel, and was read back.  No predecessor recovery request was
needed.

\section{DUPLICATION / CONTRADICTION AUDIT}

ORDERLOC1 already proves that clause order collapses a physical locator and
asks for source-matched BIND.  SYNTHESIS-0017 already imposes the general
representation-robustness gate.  ARC13 already shows that multiset equality
can be tested without sorted output in a declared finite-field stream, though
its full source work is too large for the favorable compiler target.  The
new content is the exact orbit theorem: original-order BIND is equivalent to
injective permutation encoding, its tight charge is only
$\lceil\log_2(q!)\rceil$, and canonical semantic reconstruction needs no such
tag.  No prior repository search hit stated this quantified dichotomy.

There is no contradiction with TESTS-0016's fixed-interface locator theorem,
LB-0020's conditional hereditary-pass compiler or TAGGLUE1's semantic gluing
theorem.  Actual main4/main5 still provide only conditional/partial source
constructions and explicitly leave all-input acquisition open.  Their
order-equality mechanisms are representation algorithms, not arbitrary-
decider obligations.

\section{WHAT CHANGED RELATIVE TO ROOT MAIN}

Root MAIN correctly keeps \texttt{ORDER-ROBUST-BINDING-SOURCE-PASS1}, the
model-matched forbidden realization and P versus NP \OPEN.  This run narrows
the BIND object: binding the physical order alone is insufficient and at most
linear in encoded length.  Any future pass must bind a semantic child or
opening required by exact hereditary reconstruction, survive canonicalization
and one-bit quotienting, and be acquired in the same fully charged machine.
No executive claim is strengthened; only the LSPV-4 ownership region is
reconciled.

\section{ADOPTED / REJECTED / RETIRED PRIOR COMPONENTS}

\begin{itemize}[leftmargin=1.6em]
\item \textbf{Adopted:} ORDERLOC1's clause-order counterexample;
SYNTHESIS-0017's representation gate; TAGGLUE1's demand for construction-
specific semantic BIND; LB-0020's conditional pass ledger; SAME\_REALIZATION
and full acquisition/output/replay accounting.
\item \textbf{Rejected:} physical source order, permutation count or tag
length as an unrestricted SAT-time, memory, communication or internal-state
obligation.
\item \textbf{FALSIFIED/RETIRED:}
\texttt{ORDER-BINDING-AMPLIFIES-HARDNESS1}, the unqualified repair that exact
original-order reconstruction closes the locator-to-computation gap.
\item \textbf{OPEN:} \texttt{SEMANTIC-OPENING-BIND1}, a full-language
hereditary source pass, arbitrary-decider acquisition, a matched forbidden
realization, and nonrelativizing/nonalgebrizing hinges.
\end{itemize}

\section{STALE OR OVERSTATED CLAIMS FOUND}

Any reading of ``source-matched BIND'' as necessarily preserving arbitrary
clause order is overstated: CNF semantics is permutation-invariant.  Conversely,
claiming that exact original-order reconstruction is free is also false; it
has the tight permutation-entropy cost above.  The corrected boundary is that
this cost belongs to an extra reconstruction interface and is only linear in
the native encoding.  It is not a one-bit SAT lower bound.

\section{Precise handoff}

Attempt \texttt{SEMANTIC-OPENING-BIND1}: define a total full-language
source-pass opening $y$ and a verifier such that two conflicting valid
openings for the \emph{same canonical source and child address} yield an
explicit construction-specific forbidden witness.  Prove that the opening is
needed for hereditary child reconstruction rather than source-order recovery,
and charge acquisition, failed routes, normalization, access, output, replay,
time and peak space.  Then apply LB-0020's exact exponent ledger.  Until that
arrow and a model-matched forbidden endpoint are proved, P versus NP remains
\OPEN.

\begin{thebibliography}{9}
\bibitem{MRRR12}
J.~I. Munro, R.~Raman, V.~Raman and S.~S. Rao,
``Succinct representations of permutations and functions,''
\emph{Theoretical Computer Science} 438 (2012), 74--88;
arXiv:1108.1983.

\bibitem{FLMV05}
L.~Fortnow, R.~Lipton, D.~van Melkebeek and A.~Viglas,
``Time-space lower bounds for satisfiability,''
\emph{Journal of the ACM} 52(6) (2005), 835--865,
doi:10.1145/1101821.1101822.

\bibitem{BGS75}
T.~Baker, J.~Gill and R.~Solovay,
``Relativizations of the $P=?NP$ question,''
\emph{SIAM Journal on Computing} 4(4) (1975), 431--442.

\bibitem{AW09}
S.~Aaronson and A.~Wigderson,
``Algebrization: a new barrier in complexity theory,''
\emph{ACM Transactions on Computation Theory} 1(1) (2009), Article 2.
\end{thebibliography}

\paragraph{Final status.}
\texttt{ORDER-SOURCE-BINDING-DICHOTOMY1}: \EXACT, independent review
pending.  Finite audit: \COMP.  Physical-order hardness amplification:
\RETIRED.  Semantic opening BIND, arbitrary-decider transfer, forbidden
realization, barrier escape and P versus NP: \OPEN.  P versus NP remains open.

\end{document}
% END_OF_VERIFIED_REPORT::LSPV-4::ORDERBIND1
